\documentclass[12pt]{article}
\usepackage[utf8]{inputenc}
\usepackage[T1]{fontenc}
\usepackage{hyperref}
\title{Yinhawangka Language Resources for K-12 Curriculum Development}
\author{Compiled by Claude Sonnet 5.5 (claude-sonnet-5-5) for Andr\'as Kornai}
\date{}
\begin{document}
\maketitle

\section{Introduction}
Yinhawangka (also spelled Inawangga; ISO 639-3 ywg; Glottocode yinh1234) is a Pama-Nyungan language of the Pilbara region of Western Australia. Its traditional country includes the Angelo River, Ashburton River, Hardey River, Kunderong Range, Mount Vernon Station, Rocklea and Turee Creek. Wangka Maya (the Pilbara language centre) reports that in 2004 there were estimated to be fewer than four speakers, and it classes the language as critically endangered. No children appear to be speaking it. The searches found no school curriculum, primer or textbook in the language. A curriculum would therefore have to be built from archival recordings, wordlists and community language work.

\section{Primary Educational Resources}
The searches found no published primers, readers or textbooks. Wangka Maya's plans included a dictionary and a children's book, but the searches did not confirm that either was published.
\begin{itemize}
\item Yinhawangka Language Project (community website): \url{https://www.yinhawangka.com.au/yinhawangka-language/}
\item Yinhawangka Language Project page: \url{https://www.yinhawangka.com.au/?p=55416}
\end{itemize}

\section{Linguistic and Reference Materials}
\begin{itemize}
\item Wikipedia entry on the language. It lists earlier recordings by Tindale (1957), Holmer (1963), Brand (1968) and Brown (1979). It also cites a 2008 English--Yinhawangka wordlist by Barbara Jones, published by Wangka Maya in Port Hedland, which the search did not find online. \url{https://en.wikipedia.org/wiki/Yinhawangka_language}
\item Wangka Maya, Yinhawangka language page: \url{https://www.wangkamaya.org.au/pilbara-languages/yinhawangka}
\item Glottolog reference record: \url{https://glottolog.org/resource/reference/id/181620}
\end{itemize}

\section{Religious and Community Resources}
\begin{itemize}
\item Yamatji Marlpa Aboriginal Corporation (YMAC), items tagged Yinhawangka: \url{https://www.ymac.org.au/tag/yinhawangka/}
\item Yinhawangka community site, ``Yinhawangka Discovery'': \url{https://www.yinhawangka.com.au/yinhawangka-discovery/}
\end{itemize}
No scripture translations were found.

\section{Digital and Audio Resources}
\begin{itemize}
\item Wangka Maya's language work. The centre records, transcribes and preserves Pilbara language materials and aims to produce dictionaries, stories, videos and recordings. \url{https://www.wangkamaya.org.au/about-languages/wangka-mayas-language-work}
\end{itemize}
The search did not locate specific audio files. The AIATSIS AUSTLANG database and Wangka Maya hold the historical recordings and would need to be contacted directly.

\section{Curriculum Development Context}
Wangka Maya offers a Certificate III in Aboriginal Languages, and early students included Yinhawangka speakers. With so few speakers, any school programme would be a language revival effort. It would depend on archival material, wordlists, and consultation with Yinhawangka community members and Wangka Maya. Teaching in the language for most of the curriculum is not feasible with the resources found.

\end{document}